% =====================================================================
% presentation.tex -- Beamer slide deck
% Boosting vs. Bagging: An Empirical Study of Ensemble Methods
% AI Academy -- Machine Learning Final Project, Spring 2026
% Compile: pdflatex presentation.tex (run twice for TOC)
% =====================================================================
\documentclass[aspectratio=169,11pt]{beamer}
\usetheme{Madrid}
\usecolortheme{beaver}
\usepackage[T1]{fontenc}
\usepackage{booktabs}
\usepackage{amsmath}
\usepackage{graphicx}
\usepackage{xcolor}

\setbeamertemplate{caption}[numbered]
\setbeamerfont{caption}{size=\scriptsize}
\setbeamertemplate{navigation symbols}{}

\graphicspath{{../figures/}}

\title[Boosting vs.\ Bagging]{Ensemble Methods:\\Boosting vs.\ Bagging}
\subtitle{An Empirical Study Across Four Datasets}
\author[Huseynli, Muradov, Yusifli]{Emin Huseynli \and Ziyad Muradov \and Ismayil Yusifli}
\institute[AI Academy]{AI Academy -- Machine Learning Final Project, Spring 2026}
\date[July 2026]{July 2026}

\begin{document}

% ------------------------------------------------------------------
\begin{frame}
\titlepage
\end{frame}

% ------------------------------------------------------------------
\section{Motivation}
\begin{frame}{Motivation \& Central Question}
\begin{block}{Central Question}
  \centering\textit{Under what conditions does boosting outperform bagging,\\
  and vice versa --- and \textbf{why}?}
\end{block}

\bigskip
\begin{columns}
\column{0.5\textwidth}
\textbf{Boosting (AdaBoost / SAMME)}
\begin{itemize}
  \item Re-weights hard examples sequentially
  \item Drives down bias
  \item Sensitive to label noise
\end{itemize}
\column{0.5\textwidth}
\textbf{Bagging (Random Forest)}
\begin{itemize}
  \item Bootstrap aggregation, trained independently
  \item Drives down variance
  \item More robust to noise
\end{itemize}
\end{columns}

\bigskip
\begin{itemize}
  \item All algorithms implemented from scratch in NumPy (no \texttt{sklearn} ensemble classes)
  \item Verified against \texttt{sklearn} on every dataset; 190 tests, 95\% line coverage
\end{itemize}
\end{frame}

% ------------------------------------------------------------------
\section{Methods}
\begin{frame}{Methods I -- Decision Tree \& AdaBoost}
\begin{columns}[t]
\column{0.5\textwidth}
\textbf{Decision Tree (CART)}
\begin{itemize}\small
  \item Gini impurity or entropy splitting
  \item Weighted samples (for AdaBoost); \texttt{max\_features} subsampling (for RF)
  \item Matches sklearn within 2\% on all four datasets
\end{itemize}
\column{0.5\textwidth}
\textbf{AdaBoost (SAMME)}
\begin{itemize}\small
  \item Depth-1 stumps, re-weighted each round:\\
  $\alpha_m = \ln\frac{1-\epsilon_m}{\epsilon_m} + \ln(K-1)$
  \item Misclassified samples get more weight next round
  \item Bonus: one-vs-rest wrapper, SAMME.R
\end{itemize}
\end{columns}
\end{frame}

% ------------------------------------------------------------------
\begin{frame}{Methods II -- Random Forest \& Unsupervised Pipeline}
\begin{columns}[t]
\column{0.5\textwidth}
\textbf{Random Forest}
\begin{itemize}\small
  \item $T$ trees on bootstrap samples; each split considers $\lfloor\sqrt{p}\rfloor$ random features
  \item Majority vote; out-of-bag (OOB) scoring as a free validation proxy
  \item Parallel training via \texttt{multiprocessing.Pool}
\end{itemize}
\column{0.5\textwidth}
\textbf{PCA / K-Means / DBSCAN}
\begin{itemize}\small
  \item PCA: eigendecomposition of covariance, scree + 2D projection
  \item K-Means: Lloyd's algorithm, k-means++ init, elbow method
  \item DBSCAN: density-based; $\varepsilon$ from $k$-distance knee
\end{itemize}
\end{columns}
\end{frame}

% ------------------------------------------------------------------
\section{Experimental Setup}
\begin{frame}{Datasets \& Experimental Design}
\begin{table}
\centering\small
\begin{tabular}{lrrl}
\toprule
Dataset & Samples & Features & Task \\
\midrule
Breast Cancer Wisconsin & 569 & 30 & Binary (63/37) \\
Credit Card Fraud & 284{,}807 & 30 & Binary, 0.17\% fraud \\
MNIST (3 vs.\ 8 subset) & 6{,}000 & 784 & Binary, high-dim \\
Covertype (subset) & 15{,}000 & 54 & 7-class \\
\bottomrule
\end{tabular}
\end{table}
\begin{itemize}
  \item Satisfies brief: severe imbalance ($\leq$1\%) + high-dimensional ($>$20 features)
  \item Features standardised; SMOTE applied to Credit Card Fraud training fold only
  \item \textbf{7 experiments}, seeded (\texttt{random\_state=42}), reproducible via
  \texttt{python src/experiments/run\_all.py}
\end{itemize}
\end{frame}

% ------------------------------------------------------------------
\section{Results}
\begin{frame}{Exp 1 -- Baseline Comparison}
Unpruned tree, depth-1 stump, and sklearn's tree on identical 80/20 splits.
\begin{table}
\centering\small
\begin{tabular}{lrrr}
\toprule
Dataset & Ours & sklearn & Diff.\\
\midrule
Breast Cancer & 0.9123 & 0.9123 & 0.0000 \\
Credit Card Fraud & 0.9991 & 0.9991 & 0.0000 \\
MNIST (3 vs.\ 8) & 0.9558 & 0.9583 & 0.0025 \\
Covertype & 0.7367 & 0.7370 & 0.0003 \\
\bottomrule
\end{tabular}
\end{table}
Max. gap 0.25 points (MNIST) -- well within the 2\% tolerance. A single stump
already reaches 0.921 on Breast Cancer but only 0.639 on Covertype, motivating ensembling.
\end{frame}

% ------------------------------------------------------------------
\begin{frame}{Exp 2\,\&\,3 -- Scaling: AdaBoost \& Random Forest}
\begin{columns}[t]
\column{0.5\textwidth}
\textbf{AdaBoost} ($n\_\text{estimators}: 1\to200$)
\begin{itemize}\scriptsize
  \item Training error $\to 0$ within 30--40 rounds
  \item Test error plateaus early (best round 10--150)
  \item No overfitting up to 200 rounds
\end{itemize}
\includegraphics[width=\linewidth,height=3.4cm,keepaspectratio]{experiment_2_adaboost_scaling_mnist_3_vs_8.png}
\column{0.5\textwidth}
\textbf{Random Forest} ($n\_\text{estimators}: 1\to100$)
\begin{itemize}\scriptsize
  \item Accuracy rises sharply to $\sim$20 trees, then flat
  \item OOB accuracy tracks test accuracy within 1--2 points
  \item Deeper trees help most on Covertype / MNIST
\end{itemize}
\includegraphics[width=\linewidth,height=3.4cm,keepaspectratio]{experiment_3a_rf_n_estimators_covertype.png}
\end{columns}
\end{frame}

% ------------------------------------------------------------------
\begin{frame}{Exp 4 -- Head-to-Head Comparison}
5-fold CV, $n\_\text{estimators}=100$, mean $\pm$ std accuracy:
\begin{table}
\centering\scriptsize
\begin{tabular}{lrrrr}
\toprule
Dataset & Tree & AdaBoost & Our RF & sklearn RF \\
\midrule
Breast Cancer & .910$\pm$.019 & \textbf{.967$\pm$.018} & .951$\pm$.014 & .956$\pm$.012 \\
Credit Card & .992$\pm$.001 & .995$\pm$.001 & \textbf{.998$\pm$.001} & \textbf{.998$\pm$.001} \\
MNIST (3v8) & .952$\pm$.003 & .950$\pm$.004 & .979$\pm$.003 & \textbf{.980$\pm$.002} \\
Covertype & .687$\pm$.010 & .623$\pm$.014 & .778$\pm$.014 & \textbf{.784$\pm$.006} \\
\bottomrule
\end{tabular}
\end{table}
AdaBoost wins on clean, balanced Breast Cancer; Random Forest wins (or matches
sklearn) on all three harder datasets -- widest gap on Covertype (15 points).
\end{frame}

% ------------------------------------------------------------------
\begin{frame}{Exp 5\,\&\,6 -- Noise Robustness \& Bias-Variance}
\begin{columns}[t]
\column{0.5\textwidth}
\textbf{Noise Robustness} (flip up to 20\% of labels)
\begin{itemize}\small
  \item Breast Cancer, MNIST: AdaBoost degrades 3--6$\times$ faster than RF
  \item Credit Card, Covertype: pattern reverses, but gap $<$1.6 points
\end{itemize}
\column{0.5\textwidth}
\textbf{Bias--Variance} ($B{=}100$, Breast Cancer)
\begin{tabular}{lrr}
\toprule
Model & Bias$^2$ & Variance \\
\midrule
Single Tree & 0.0702 & 0.0626 \\
AdaBoost & 0.0351 & 0.0257 \\
Random Forest & 0.0585 & 0.0146 \\
\bottomrule
\end{tabular}
\end{columns}
\vfill
\small AdaBoost halves bias$^2$; Random Forest quarters variance -- the textbook mechanism.
\end{frame}

% ------------------------------------------------------------------
\begin{frame}{Exp 7 \& Bonus -- Unsupervised Analysis, SAMME.R vs.\ OvR}
\begin{columns}[t]
\column{0.5\textwidth}
\textbf{Unsupervised (PCA / K-Means / DBSCAN)}
\begin{itemize}\scriptsize
  \item Breast Cancer: K-Means ARI 0.654 -- classes well separated
  \item Credit Card Fraud: K-Means ARI $\approx$0 -- fraud invisible to distance clustering
  \item Covertype: DBSCAN (0.160) beats K-Means (0.044) -- non-convex geometry
\end{itemize}
\column{0.5\textwidth}
\textbf{SAMME.R vs.\ OvR} (Covertype, $\sim$100 stumps)
\begin{itemize}\scriptsize
  \item SAMME.R: 0.674 acc., log-loss 2.650, 4.62s
  \item OvR (compute-matched): 0.678 acc., 3.93s -- essentially tied
  \item Naive OvR (7$\times$ stumps): +2 acc.\ points for 7$\times$ the time -- poor trade-off
\end{itemize}
\end{columns}
\end{frame}

% ------------------------------------------------------------------
\section{Discussion \& Conclusion}
\begin{frame}{Key Findings \& Conclusion}
\begin{itemize}\small
  \item \textbf{Boosting wins:} binary, reasonably separable classes, clean labels --
  lowest bias$^2$, best accuracy on Breast Cancer
  \item \textbf{Bagging wins:} severe class imbalance, high dimensionality, or $>$2
  classes -- lowest variance, native multi-class handling
  \item Implemented Decision Tree, AdaBoost (SAMME), Random Forest, PCA, K-Means,
  DBSCAN from scratch (NumPy only); verified vs.\ sklearn within 2\%; 190 tests,
  95\% line coverage
  \item \textbf{Limitations:} RF parallelism untested at scale; noise-robustness
  reversal on 2 datasets needs deeper study
  \item \textbf{Future work:} from-scratch Gradient Boosting Machine; extend
  bias-variance decomposition to all four datasets
\end{itemize}
\vfill
\begin{center}
\Large\textbf{Thank you -- Questions?}
\end{center}
\end{frame}

\end{document}
